%% appendix-summaries.tex -- proofs for the sections on pushdown summaries
%% (sec:sm), condition (U) (sec:un), and two-word holes with self events
%% (sec:ex): summaries and the lowest point (by induction along a derivation),
%% duality, the least closure, deadness from states, the deadness theorem, and
%% the reduction of (II) to holes of the extracted instances. The proofs for the
%% smaller unit reader of sec:un-smaller are in appendix-complexity.

\section{Proofs for Summaries, Condition (U), and Extracted Instances}\label{app:sm}

\subsection{Summaries and the Lowest Point}\label{app:sm-summaries}

Throughout, $\chi$ is a system over $V$, pure derivations are those of
$\chi$, that is, of the copy and constructor literals of $\chi^\pm$
(\cref{sec:sm-lowest}), and the automata are those of \cref{def:sm-automata}.

We prove \cref{lem:sm-summaries,lem:sm-lowest-point} together: first the
derivations given by the summaries, then the lowest point by induction along a
derivation, and finally the converse directions of \cref{lem:sm-summaries} as
the special cases of the lowest point with an empty word. Recall that every
set $P^\chi_p(w)$ is upward closed and every set $Q^\chi_q(w)$ downward closed
under $D_\chi$.

\emph{Step 1: summaries give derivations.} We show the forward directions of
\cref{lem:sm-summaries}(a)--(c), lifted to an arbitrary suffix $\sigma$. (a)
Induct on the derivation of $D_\chi(a,b)$ by the rules of
\cref{def:sm-automata}, constructing a pure derivation from $(a,\sigma)$ to
$(b,\sigma)$ for every $\sigma$ at once. Reflexivity gives the empty
derivation, a copy literal $a\le b$ of $\chi^\pm$ one copy, and transitivity a
concatenation. For the decomposition rule, which concludes $D_\chi(a_i,b_i)$
from $f(a_j)_j\le u$, $D_\chi(u,v)$, and $v\le f(b_j)_j$: push $i$ from
$(a_i,\sigma)$ to $(u,i\sigma)$, follow the derivation from $(u,i\sigma)$ to
$(v,i\sigma)$ given by the induction hypothesis, and pop $i$ to $(b_i,\sigma)$.
(b) Induct on $\gamma$. For $\gamma=\eps$, $P^\chi_p(\eps)=\upc\chi p$, and (a)
applies. If $b\in P^\chi_p(\gamma i)$, there are $r\in P^\chi_p(\gamma)$ and a
literal $r\le f(b_j)_j$ of $\chi^\pm$ with $D_\chi(b_i,b)$; the induction
hypothesis with $i\sigma$ in place of $\sigma$, a pop of $i$, and (a) compose to
a derivation from $(p,\gamma i\sigma)$ to $(b,\sigma)$. (c) Symmetrically, if
$a\in Q^\chi_q(\beta i)$, there are $r\in Q^\chi_q(\beta)$ and a literal
$f(a_j)_j\le r$ with $D_\chi(a,a_i)$; (a), a push of $i$ from $(a_i,\sigma)$ to
$(r,i\sigma)$, and the induction hypothesis compose.

\emph{Step 2: three facts on the lower automata.} For signed variables $q,q'$,
words $\beta$ and letters $i$:
\begin{enumerate}[label=(L\arabic*),leftmargin=*]
\item if $D_\chi(q',q)$, then $Q^\chi_{q'}(\beta)\subseteq Q^\chi_q(\beta)$;
\item if $f(a_j)_j\le q$ is a literal of $\chi^\pm$ with $a_i=q'$, then
  $Q^\chi_{q'}(\beta)\subseteq Q^\chi_q(i\beta)$;
\item if $q'\le f(b_j)_j$ is a literal of $\chi^\pm$, then
  $Q^\chi_{q'}(i\beta)\subseteq Q^\chi_{b_i}(\beta)$.
\end{enumerate}
All three follow by induction on $\beta$, since $\qstep j$ is monotone. For
$\beta=\eps$: (L1) is transitivity of $D_\chi$; in (L2), an element $a$ of
$\downc\chi{q'}$ lies in $\qstep i(\downc\chi q)$, witnessed by $q$ itself and
the literal $f(a_j)_j\le q$, since $D_\chi(a,a_i)$; in (L3), an element $a$ of
$\qstep i(\downc\chi{q'})$ comes with $r$, $D_\chi(r,q')$, a literal
$f(a'_j)_j\le r$ and $D_\chi(a,a'_i)$, and the decomposition rule applied to
$f(a'_j)_j\le r$, $D_\chi(r,q')$ and $q'\le f(b_j)_j$ gives
$D_\chi(a'_i,b_i)$, hence $a\in\downc\chi{b_i}$.

\begin{proof}[Proof of \cref{lem:sm-lowest-point}]
If $\wsuf wi=\wsuf ve=\sigma$ and $r\in P^\chi_{p,i}(w)\cap Q^\chi_{q,e}(v)$,
then Step~1 gives pure derivations
$(p,w)=(p,\wpre wi\sigma)\to^*_\chi(r,\sigma)\to^*_\chi(q,\wpre ve\sigma)=(q,v)$.

Conversely, consider a pure derivation from $(p,w)$. We show by induction on
its length that every atom $(q',v')$ on it admits cuts $i\le\len w$ and
$e\le\len{v'}$ with $\wsuf wi=\wsuf{v'}e$ and
$P^\chi_{p,i}(w)\cap Q^\chi_{q',e}(v')\ne\varnothing$; for the last atom this
is the claim. For the first atom take $i=e=0$: $p\in\upc\chi p\cap\downc\chi p$.
Let the invariant hold at $(q',v')$ with cuts $i,e$ and a signed variable $r$
in the intersection, and let the next move lead to $(q'',v'')$.
\begin{itemize}[leftmargin=*]
\item A copy $q'\le q''$: then $v''=v'$, and $r\in Q^\chi_{q'',e}(v'')$ by
  (L1); keep $i$ and $e$.
\item A push of $j$ by a literal $f(a_k)_k\le q''$ with $a_j=q'$: then
  $v''=jv'$, and $r\in Q^\chi_{q''}(j\,\wpre{v'}e)=Q^\chi_{q'',e+1}(v'')$ by
  (L2); keep $i$ and take $e+1$.
\item A pop of $j$ by a literal $q'\le f(b_k)_k$: then $v'=jv''$ and
  $q''=b_j$. If $e\ge1$, then $\wpre{v'}e=j\,\wpre{v''}{e-1}$, and
  $r\in Q^\chi_{b_j,e-1}(v'')$ by (L3); keep $i$ and take $e-1$. If $e=0$,
  then $r\in\downc\chi{q'}$, so $q'\in P^\chi_{p,i}(w)$ by upward closure, and
  $\wsuf wi=v'=jv''$; hence $i<\len w$, $w[i]=j$, and
  $b_j\in\pstep j(P^\chi_{p,i}(w))=P^\chi_{p,i+1}(w)$, witnessed by $q'$ and
  the literal. Take $i+1$ and $e=0$, with the signed variable
  $b_j\in\downc\chi{b_j}$.
\end{itemize}
In every case $\wsuf w{i}=\wsuf{v''}{e}$ for the new cuts.
\end{proof}

\begin{proof}[Proof of \cref{lem:sm-summaries}]
The forward directions are Step~1 with $\sigma=\eps$. For the converse of (b),
apply \cref{lem:sm-lowest-point} to a pure derivation from $(p,\gamma)$ to
$(b,\eps)$: the only cut of $\eps$ is $0$, so $\wsuf\gamma i=\eps$ forces
$i=\len\gamma$, and some $r\in P^\chi_p(\gamma)\cap\downc\chi b$ gives
$b\in P^\chi_p(\gamma)$ by upward closure. The converse of (c) is symmetric,
by downward closure, and that of (a) is the case $\gamma=\eps$ of (b), since
$P^\chi_a(\eps)=\upc\chi a$.
\end{proof}

\begin{proof}[Proof of \cref{lem:sm-duality}]
By \cref{lem:lf-duality}, reversing a pure derivation and flipping every
signed variable gives a pure derivation, and the dual of an atom has the same
word. (1) By \cref{lem:sm-summaries}(a), $D_\chi(a,b)$ iff
$(a,\eps)\to^*_\chi(b,\eps)$, iff $(\bar b,\eps)\to^*_\chi(\bar a,\eps)$, iff
$D_\chi(\bar b,\bar a)$. Hence $r\in\downc\chi{\bar q}$ iff $D_\chi(r,\bar q)$
iff $D_\chi(q,\bar r)$ iff $r\in\overline{\upc\chi q}$. (2) By
\cref{lem:sm-summaries}(c,b), $r\in Q^\chi_{\bar q}(w)$ iff
$(r,\eps)\to^*_\chi(\bar q,w)$, iff $(q,w)\to^*_\chi(\bar r,\eps)$, iff
$\bar r\in P^\chi_q(w)$. (3) By \cref{lem:sm-summaries}, $D_\chi(a,b)$ gives
$(a,\eps)\to^*_\chi(b,\eps)$, $r\in P^\chi_p(w)$ gives
$(p,w)\to^*_\chi(r,\eps)$, and $r\in Q^\chi_q(w)$ gives
$(r,\eps)\to^*_\chi(q,w)$. The sign invariant of \cref{lem:lf-moves}(2),
extended to pure derivations, gives the three claims, since
$\parity(\eps)=0$.
\end{proof}

\subsection{Least Closures and Deadness}\label{app:sm-dead}

\begin{proof}[Proof of \cref{lem:sm-mu}]
By definition, $(q,v)\in\mu_\chi$ iff $(q,v)\in\Reach_\chi(p,\eps)$ for some
fact $(p,\eps)$, iff, by \cref{lem:lf-ext-first},
$(p,\alpha)\to^*_\chi(q,v)$ for some $\alpha$. By \cref{lem:sm-lowest-point},
this holds iff there are cuts $i\le\len\alpha$ and $e\le\len v$ with
$\wsuf\alpha i=\wsuf ve$ and
$P^\chi_p(\wpre\alpha i)\cap Q^\chi_{q,e}(v)\ne\varnothing$. If so, then
$Q^\chi_{q,e}(v)$ meets $P^\chi_p(\gamma)\subseteq\Fup\chi$ for
$\gamma=\wpre\alpha i$. Conversely, if $r\in P^\chi_p(\gamma)\cap Q^\chi_{q,e}(v)$
for a fact $(p,\eps)$, a word $\gamma$, and a cut $e$, take
$\alpha=\gamma\,\wsuf ve$ and $i=\len\gamma$.

For the bound, if $r\in P^\chi_e(\gamma)$, then the state $P^\chi_e(\gamma)$ of
the upper automaton of $e$, which has $4^{\len V}$ states, is reached by a word
that visits no state twice, of length less than $4^{\len V}$.
\end{proof}

\begin{proof}[Proof of \cref{lem:sm-dead-states}]
We first show that
\begin{equation}\label{eq:sm-union}
 P^\chi_p(wz)=\bigcup_{r\in P^\chi_p(w)}P^\chi_r(z)\qquad\text{for all words }w,z.
\end{equation}
For $z=\eps$ this holds because $P^\chi_p(w)$ is upward closed under $D_\chi$
and $r\in\upc\chi r=P^\chi_r(\eps)$. The induction step from $z$ to $zi$ applies
$\pstep i$, which distributes over unions.

(1) By \cref{lem:sm-roots}(3) and \cref{eq:sm-union}, $\DeadU_\chi(p,w)$ iff
some $r\in P^\chi_p(w)$ and some $z$ have
$P^\chi_r(z)\cap\Ndown\chi\ne\varnothing$, that is, by \cref{lem:sm-roots}(3)
for the atom $(r,\eps)$, iff some $r\in P^\chi_p(w)$ has
$\DeadU_\chi(r,\eps)$.

(2) This is \cref{lem:sm-mu} for the atom $(\bar p,w)$, together with
$Q^\chi_{\bar p}=\overline{P^\chi_p}$ (\cref{lem:sm-duality}(2)).

(3) Suppose $\DeadT_\chi(p,w)$. If $\DeadU_\chi(p,w)$, then (1) gives some
$r\in P^\chi_p(w)$ with $\DeadU_\chi(r,\eps)$, hence $r\in\DeadT^\eps_\chi$.
Otherwise $(\bar p,w\delta)\in\mu_\chi$ for some $\delta$, and by (2) applied to
$w\delta$ there is a cut $e\le\len{w\delta}$ with
$\overline{P^\chi_p(\wpre{(w\delta)}e)}\cap\Fup\chi\ne\varnothing$. If
$e\le\len w$, then $(\bar p,w)\in\mu_\chi$ by (2). If $e>\len w$, write
$\wpre{(w\delta)}e=wz$; by \cref{eq:sm-union}, some $r\in P^\chi_p(w)$ has
$\overline{P^\chi_r(z)}\cap\Fup\chi\ne\varnothing$, and (2) for the atom
$(\bar r,z)$ gives $(\bar r,z)\in\mu_\chi$, so $r\in\DeadT^\eps_\chi$.

Conversely, $(\bar p,w)\in\mu_\chi$ gives $\DeadT_\chi(p,w)$ with $\delta=\eps$.
Let $r\in P^\chi_p(w)\cap\DeadT^\eps_\chi$; then $(p,w)\to^*_\chi(r,\eps)$ by
\cref{lem:sm-summaries}(b). If $\DeadU_\chi(r,\eps)$, then
$\Reach_\chi(r,\eps)\subseteq\Reach_\chi(p,w)$ contains a forbidden root. If
$(\bar r,\delta)\in\mu_\chi$, then duality reverses the derivation to
$(\bar r,\eps)\to^*_\chi(\bar p,w)$, \cref{lem:lf-moves}(1) lifts it to
$(\bar r,\delta)\to^*_\chi(\bar p,w\delta)$, and $(\bar p,w\delta)\in\mu_\chi$
because $\mu_\chi$ is closed under the rules (\cref{lem:lf-closure}(1)).

(4) The first part of $\DeadT_\chi$ is inherited by prefixes by
\cref{lem:sm-roots}(4), and $(\bar p,wv\delta)=(\bar p,w(v\delta))$.

(5) By \cref{lem:sm-roots}(3), $r\in\DeadU^\eps_\chi$ iff some word $z$ has
$P^\chi_r(z)\cap\Ndown\chi\ne\varnothing$. This is the reachability of a set of
states of the upper automaton of $r$, which has $4^{\len V}$ states, so words
$z$ of length less than $4^{\len V}$ suffice. Similarly, by (2) and
\cref{eq:sm-union}, the set $M$ of signed variables $r$ with
$(\bar r,\delta)\in\mu_\chi$ for some $\delta$ consists of the $r$ with
$P^\chi_r(z)\cap\overline{\Fup\chi}\ne\varnothing$ for some word $z$, and the
same bound applies. Finally $\DeadT^\eps_\chi=\DeadU^\eps_\chi\cup M$, and
$\Fup\chi$ is computable by \cref{lem:sm-mu}.
\end{proof}

\begin{proof}[Proof of \cref{thm:sm-dead}]
By \cref{cor:lf-dead}, $(p,w)$ is dead iff (a) $\DeadU_\chi(p,w)$, or (b)
$(\bar p,w\delta)\in\mu_\chi$ for some $\delta$, or (c) $(p,w)$ has a
self-collision, that is, $(p,w\alpha)\to^*_\chi(\bar p,w\delta)$ for some
words $\alpha$ and $\delta$. Causes (a) and (b) together are
$\DeadT_\chi(p,w)$, and the causes may overlap. It remains to show that a
self-collision implies $\DeadT_\chi(p,w)$ or a self event, and that a self
event implies a self-collision.

Assume a self-collision. By \cref{lem:sm-lowest-point,lem:sm-duality}(2), there are cuts
$i\le\len{w\alpha}$ and $j\le\len{w\delta}$ with
$\wsuf{(w\alpha)}i=\wsuf{(w\delta)}j$ and
\[
 P^\chi_p(\wpre{(w\alpha)}i)\cap\overline{P^\chi_p(\wpre{(w\delta)}j)}
 \ne\varnothing .
\]
In particular both sets $P^\chi_p(\wpre{(w\alpha)}i)$ and
$P^\chi_p(\wpre{(w\delta)}j)$ are nonempty. If $i>\len w$, then
$\wpre{(w\alpha)}i$ is a proper extension of $w$ with a nonempty upper set, and
the pop bound (\cref{lem:sm-roots}(2) with $k=\len w$) gives $\DeadU_\chi(p,w)$; the
case $j>\len w$ is the same. Otherwise $i,j\le\len w$, and the equation reads
$\wsuf wi\,\alpha=\wsuf wj\,\delta$, so $\wsuf wi$ and $\wsuf wj$ are prefixes
of one word and hence comparable in the prefix order, and the displayed
condition reads $P^\chi_{p,i}(w)\cap\overline{P^\chi_{p,j}(w)}\ne\varnothing$. By
sign coherence (\cref{lem:sm-duality}(3)), the elements of $P^\chi_{p,i}(w)$
have the sign $\sgn(p)\oplus\parity(\wpre wi)$ and those of
$\overline{P^\chi_{p,i}(w)}$ the opposite sign, so $i\ne j$. If $i<j$, then
$\wsuf wj$ is the shorter suffix, so $\wsuf wj\pre\wsuf wi$, and $(i,j)$ is a
self event. If $j<i$, flipping both sets gives
$P^\chi_{p,j}(w)\cap\overline{P^\chi_{p,i}(w)}\ne\varnothing$ and
$\wsuf wi\pre\wsuf wj$, so $(j,i)$ is a self event.

Conversely, let $(i,j)$ be a self event of $(p,w)$. Then $\wsuf wi=\wsuf wj\,\delta$
for some word $\delta$, and the guard says, by \cref{lem:sm-duality}(2), that
$P^\chi_{p,i}(w)\cap Q^\chi_{\bar p}(\wpre wj)\ne\varnothing$. Since
$\wpre{(w\delta)}j=\wpre wj$ and $\wsuf{(w\delta)}j=\wsuf wj\,\delta=\wsuf wi$,
\cref{lem:sm-lowest-point} gives $(p,w)\to^*_\chi(\bar p,w\delta)$, a
self-collision with $\alpha=\eps$.
\end{proof}

\subsection{Extracted Instances}\label{app:sm-extraction}

\begin{proof}[Proof of \cref{thm:ex-reduction}]
Write $\zeta_e$ for the state of $B=\Bsrc_p$ at the cut $e$ of $t$. By
\cref{def:ex-instance,lem:un-source}, $A_s(u)=Q^\varphi_{q,s}(u)$,
$A'_s(u)=Q^\psi_{q,s}(u)$, $B'_e(t)=P^\psi_{p,e}(t)$, and
$\zeta_e=(P^\varphi_{p,e}(t),[(\bar p,\wpre te)\in\mu_\varphi],P^\psi_{p,e}(t))$.
We translate the conditions of \cref{def:ex-ss2self}.
\begin{itemize}[leftmargin=*]
\item Condition (a) says that $Q^\varphi_{q,s}(u)\cap\Fup\varphi=\varnothing$
  for every cut $s$, that is, $(q,u)\notin\mu_\varphi$ (\cref{lem:sm-mu}).
\item By \cref{lem:un-source}(2), a cut that is rejected stays rejected at
  every later cut, so condition (b) says that the last cut of $t$ is rejected:
  neither $\DeadT_\varphi(p,t)$ nor $\DeadU_\psi(p,t)$ holds.
\item Condition (d) says that some cuts $s\le\len u$ and $e\le\len t$ have
  $\wsuf te=\wsuf us$ and $P^\psi_{p,e}(t)\cap Q^\psi_{q,s}(u)\ne\varnothing$;
  by \cref{lem:sm-lowest-point} for $\psi$, this says $(p,t)\to^*_\psi(q,u)$.
\item Assume (b). At the cut $J(t)$, $\mathrm{dt}$ fails, so
  $\DeadU_\varphi(p,\wpre t{J(t)})$ fails, and the pop bound (\cref{lem:sm-roots}(2))
  gives $P^\varphi_{p,e}(t)=\varnothing$ for every $e>J(t)$. A guard
  $Q^\varphi_{q,s}(u)\cap P^\varphi_{p,e}(t)\ne\varnothing$ or a self guard
  $P^\varphi_{p,i}(t)\cap\overline{P^\varphi_{p,j}(t)}\ne\varnothing$ therefore
  forces $e\le J(t)$, respectively $j\le J(t)$. Hence, under (b), condition (c)
  says that there are no cuts $e\le\len t$ and $s\le\len u$ with
  $\wsuf te\pre\wsuf us$ and
  $P^\varphi_{p,e}(t)\cap Q^\varphi_{q,s}(u)\ne\varnothing$, and condition (e)
  says that $(p,t)$ has no self event in $\varphi$
  (\cref{lem:un-source}(3)).
\end{itemize}

Let $(u,t)$ be a hole. By (b) and (e), $(p,t)$ is not dead by support and has no
self event, so it is live by \cref{thm:sm-dead}, since $\varphi$ is
satisfiable. By (d), $(p,t)\to^*_\psi(q,u)$. By (a) and (c), and by
\cref{lem:sm-membership} for the live seed $(p,t)$ and the atom $(q,u)$,
$(q,u)\notin\Cl_\varphi(p,t)$.

Conversely, let $(p,t)$ be live in $\varphi$, $(p,t)\to^*_\psi(q,u)$, and
$(q,u)\notin\Cl_\varphi(p,t)$. By \cref{thm:sm-dead}, $\DeadT_\varphi(p,t)$
fails and $(p,t)$ has no self event in $\varphi$; by \CondU\ and
\cref{lem:un-dead}, $\DeadU_\psi(p,t)$ fails. So (b) and (e) hold. Condition
(d) holds by the derivation. Since $(q,u)\notin\Cl_\varphi(p,t)$,
\cref{lem:sm-membership} gives $(q,u)\notin\mu_\varphi$, which is (a), and the
absence of the cuts described above, which is (c).

(1) By \cref{lem:ch-pure}(2), \CondII\ fails iff some visible atom $(p,t)$
that is live in $\varphi$ and some visible atom $(q,u)$ satisfy
$(p,t)\to^*_\psi(q,u)$ and $(q,u)\notin\Cl_\varphi(p,t)$. By the equivalence
just proved, this happens iff some $I_{pq}$ has a hole.

(2) Let $(u,t)$ be a hole of $I_{pq}$. Then $(p,t)$ is live, so
$S=\Cl_\varphi(p,t)$ is a model of $\varphi$ (\cref{lem:lf-closure}(2)). If
$S=T|_X$ for a model $T$ of $\psi$, then $T$ contains $(p,t)$ and is closed
under the rules of $\psi$, so it contains $(q,u)$, which is visible; hence
$(q,u)\in S$, a contradiction. By \cref{lem:lf-clauses}, the solution of
$\varphi$ whose set of atoms is $S$ has no extension to a solution of $\psi$.
\end{proof}
